\documentclass[11pt]{article}

\usepackage[margin=1in]{geometry}
\usepackage{amsmath,amssymb,amsthm,mathtools}
\usepackage{booktabs}
\usepackage{enumitem}
\usepackage[hidelinks]{hyperref}

\newtheorem{theorem}{Theorem}
\newtheorem{proposition}{Proposition}
\newtheorem{lemma}{Lemma}
\newtheorem{corollary}{Corollary}
\theoremstyle{definition}
\newtheorem{definition}{Definition}
\newtheorem{assumption}{Assumption}
\newtheorem{remark}{Remark}
\newtheorem{example}{Example}

\newcommand{\R}{\mathbb{R}}
\newcommand{\Rnn}{\mathbb{R}_{\ge 0}}
\newcommand{\Af}{A_{\mathrm{fusion}}}
\newcommand{\Smat}{S}
\newcommand{\Wmat}{W}
\newcommand{\Gmat}{G}
\newcommand{\Shat}{\widehat S}
\newcommand{\What}{\widehat W}
\newcommand{\Bpchat}{\widehat B^{PC}}
\newcommand{\Bcphat}{\widehat B^{CP}}
\newcommand{\diag}{\operatorname{diag}}
\newcommand{\norm}[1]{\left\lVert #1 \right\rVert}

\title{A Formal Analysis of the SENA Fusion Matrix:\\
A Typed Supra-Adjacency Model for Social-Epistemic Network Analysis}

\author{Draft mathematical note for the SENA project}
\date{June 8, 2026; directed bridge contract revised July 10, 2026}

\begin{document}

\maketitle

\begin{abstract}
Social Network Analysis (SNA) represents relations among persons, while
Epistemic Network Analysis (ENA) represents relations among coded epistemic
elements. The SENA fusion formula combines a person-person social matrix, a
code-code epistemic co-occurrence matrix, and typed person-to-code and
code-to-person bridge matrices into one block matrix. This paper gives a formal mathematical analysis
of that formula. We show that, under explicit dimensional, nonnegativity,
symmetry, and normalization assumptions, the SENA fusion matrix is a valid
weighted typed adjacency matrix on a heterogeneous social-epistemic graph. We
prove dimensional consistency, symmetry conditions, Laplacian positive
semidefiniteness, coherent boundary cases, and the distinction between a joint
adjacency model and an ENA projection. We also identify failure modes and
statistical validation requirements. The main conclusion is that SENA should be
interpreted not as a simple visual overlay of SNA and ENA, but as a typed
supra-adjacency representation of a social-epistemic nexus. Its novelty is
domain-specific formalization and validation, not the invention of block
adjacency matrices themselves.
\end{abstract}

\noindent\textbf{Keywords:} social network analysis; epistemic network
analysis; heterogeneous networks; multilayer networks; supra-adjacency matrix;
collaborative learning analytics; SENA.

\noindent\textbf{Mathematics Subject Classification:} 05C50, 05C82, 62H25,
91D30.

\section{Introduction}

Collaborative learning is simultaneously social and epistemic. Learners
interact with particular people while producing, connecting, challenging, and
refining particular ideas. SNA gives mature tools for person-person relations,
such as degree, density, centrality, communities, and brokerage
\cite{wasserman1994}. ENA gives a
formal way to model co-occurrence relations among coded epistemic elements and
to project units of analysis into a low-dimensional epistemic space
\cite{shaffer2016,bowman2021}.

The methodological gap is not that SNA and ENA cannot both be applied to the
same data. Prior work such as SENS and iSENS already argues for combining
social and epistemic perspectives \cite{gasevic2019,swiecki2020}. The sharper mathematical question is whether
one can construct a single object that contains person-person, code-code, and
person-code relations while preserving their type distinctions and while
remaining compatible with standard network analysis.

The proposed SENA formula is
\begin{equation}
\label{eq:fusion}
\Af =
\begin{bmatrix}
\alpha \Shat & \gamma \Bpchat \\
\gamma \Bcphat & \beta \What
\end{bmatrix},
\end{equation}
where \(\Shat\) is a normalized social matrix, \(\What\) is a normalized
epistemic co-occurrence matrix, \(\Bpchat\) is a normalized person-to-code
bridge, \(\Bcphat\) is a normalized code-to-person bridge, and
\(\alpha,\beta,\gamma\ge 0\) are layer weights. The implementation retains
\(B\) as a compatibility alias for \(B^{PC}\). Without independent
code-to-person evidence, \(B^{CP}=(B^{PC})^\top\); otherwise the bridge blocks
are estimated independently and the fusion matrix is directed.

This paper analyzes the fusion matrix as a mathematical object. The
central claim is:

\begin{quote}
The SENA formula is mathematically valid as a normalized weighted heterogeneous
adjacency matrix on the typed node set \(P\cup C\). It is not, by itself, an
ENA projection, a causal model, or an inferential test.
\end{quote}

\section{Related Mathematical and Methodological Context}

The formula sits at the intersection of four established literatures. First,
SNA supplies the person-person graph layer \cite{wasserman1994}. Second, ENA
supplies the code-code co-occurrence logic and, separately, projection methods
for units of analysis \cite{shaffer2016,siebert2017,bowman2021}. Third,
multilayer network theory supplies supra-adjacency representations for systems
with multiple layers or relation types \cite{kivela2014,mucha2010}. Fourth,
heterogeneous information network theory supplies typed nodes and typed edges
\cite{sun2012}.

SENS and iSENS are the nearest learning-analytics precedents. They establish
that SNA and ENA can provide complementary information about collaborative
learning roles, discourse, group differences, and performance
\cite{gasevic2019,swiecki2020}. However, their main contribution is an
integrated analytic strategy rather than the specific block object represented
by the SENA fusion matrix. Conversely, multilayer and heterogeneous-network
theory make this block construction mathematically unsurprising, but they do
not supply the learning-analytics semantics of \(S\), \(W\), \(B^{PC}\),
\(B^{CP}\), and \(G\).

Consequently, a calibrated novelty claim is required. SENA does not invent
block matrices or multilayer graphs. Its contribution is the domain-specific
formalization of a social-epistemic graph that keeps persons and epistemic
codes as different node types while retaining social ties, epistemic
co-occurrences, and contribution bridges in one traceable model object.

\section{Notation and Empirical Matrices}

Let
\[
P=\{p_1,\ldots,p_n\},\qquad C=\{c_1,\ldots,c_m\}
\]
denote the person set and the code set. The typed node set is
\[
V=P\cup C,\qquad |V|=n+m,\qquad P\cap C=\varnothing.
\]

\begin{definition}[Social layer]
The social layer is a matrix
\[
\Smat\in \Rnn^{n\times n},
\]
where \(S_{ij}\) denotes the observed interaction strength from \(p_i\) to
\(p_j\). In an undirected model, \(S_{ij}=S_{ji}\).
\end{definition}

\begin{definition}[Epistemic co-occurrence layer]
The epistemic layer is a matrix
\[
\Wmat\in \Rnn^{m\times m},
\]
where \(W_{ab}\) denotes the co-occurrence strength between codes \(c_a\) and
\(c_b\) within a specified unit, stanza, or moving window. In the usual ENA
co-occurrence construction, \(W_{ab}=W_{ba}\) and diagonal entries are either
zero or treated separately as code frequencies.
\end{definition}

\begin{definition}[Typed bridge layers]
The person-to-code bridge is
\[
B^{PC}\in \Rnn^{n\times m},
\]
where \(B^{PC}_{ia}\) denotes the contribution, use, or activation strength of
code \(c_a\) by person \(p_i\). The code-to-person bridge is
\[
B^{CP}\in \Rnn^{m\times n},
\]
where \(B^{CP}_{ai}\) denotes explicit target, uptake, feedback, or
recommendation evidence from code \(c_a\) toward person \(p_i\). The symbol
\(B\) remains a compatibility alias for \(B^{PC}\).
\end{definition}

\begin{definition}[Layer normalization]
A layer normalization is a map \(N_L:M_L\mapsto \widehat M_L\) applied
separately to \(L\in\{S,W,B^{PC},B^{CP}\}\). Examples include
\[
\widehat M=\frac{M}{\max_{uv}|M_{uv}|},\qquad
\widehat M=\frac{\log(1+M)}{\max_{uv}\log(1+M_{uv})},\qquad
\widehat M=\frac{M}{\norm{M}_F}.
\]
\end{definition}

\section{A Generative Construction}

Let \(T=\{t_1,\ldots,t_q\}\) be a set of discourse windows. Let
\[
X_{ta}=\text{activation strength of code }c_a\text{ in window }t,
\]
\[
Y_{it}=\text{participation strength of person }p_i\text{ in window }t,
\]
and let \(R_{ij}\) be the observed social interaction strength from \(p_i\) to
\(p_j\). A common SENA-compatible construction is
\begin{equation}
S_{ij}=R_{ij} \quad\text{or}\quad S_{ij}=R_{ij}+R_{ji}.
\end{equation}
\begin{equation}
W_{ab}=\sum_{t=1}^q X_{ta}X_{tb}, \quad a\ne b.
\end{equation}
\begin{equation}
B^{PC}_{ia}=\sum_{t=1}^q Y_{it}X_{ta}.
\end{equation}
The second equation is an ENA-style code co-occurrence object before
dimensional reduction. The third equation is an affiliation or contribution
object connecting persons to codes. For coded segment \(r\) with code weight
\(x_{ra}\) and explicitly declared target-person set \(T(r)\), the independent
reverse bridge is estimated by
\begin{equation}
B^{CP}_{ai}=\sum_r x_{ra}\mathbf 1\{p_i\in T(r)\}.
\end{equation}
When no valid target-person evidence exists, the contract requires the fallback
\(B^{CP}=(B^{PC})^\top\).

\section{Assumptions}

\begin{assumption}[Dimensional compatibility]
\(\Smat\in\R^{n\times n}\), \(\Wmat\in\R^{m\times m}\), and
\(B^{PC}\in\R^{n\times m}\), \(B^{CP}\in\R^{m\times n}\).
\end{assumption}

\begin{assumption}[Nonnegativity]
\(\Smat,\Wmat,B^{PC},B^{CP}\) have nonnegative entries and
\(\alpha,\beta,\gamma\ge 0\).
\end{assumption}

\begin{assumption}[Undirected baseline]
For the undirected SENA baseline, \(\Smat=\Smat^\top\) and
\(\Wmat=\Wmat^\top\), and \(B^{CP}=(B^{PC})^\top\).
\end{assumption}

\begin{assumption}[Explicit normalization]
The matrices entering the fusion matrix are normalized within layers
before cross-layer scaling by \(\alpha,\beta,\gamma\).
\end{assumption}

\begin{assumption}[Typed-node semantics]
The labels \(p_i\) and \(c_a\) belong to different node types. Distances or
centralities computed from \(\Af\) must be interpreted with respect to this
typed graph, not as if all nodes were semantically identical.
\end{assumption}

\section{Main Results}

\begin{proposition}[Dimensional consistency]
Under Assumption 1, the matrix \(\Af\) is an
\((n+m)\times(n+m)\) matrix.
\end{proposition}

\begin{proof}
The upper-left block has \(n\) rows and \(n\) columns. The upper-right block has
\(n\) rows and \(m\) columns. The lower-left block has \(m\) rows and \(n\)
columns. The lower-right block has \(m\) rows and \(m\) columns. Therefore the
block matrix has \(n+m\) rows and \(n+m\) columns.
\end{proof}

\begin{theorem}[Weighted undirected heterogeneous graph]
Under Assumptions 1--3, \(\Af\) is a symmetric nonnegative adjacency matrix on
the typed node set \(V=P\cup C\).
\end{theorem}

\begin{proof}
Nonnegativity follows from nonnegative entries and nonnegative layer weights.
For symmetry,
\[
\Af^\top =
\begin{bmatrix}
(\alpha\Shat)^\top & (\gamma\Bcphat)^\top\\
(\gamma\Bpchat)^\top & (\beta\What)^\top
\end{bmatrix}
=
\begin{bmatrix}
\alpha\Shat & \gamma\Bpchat\\
\gamma\Bcphat & \beta\What
\end{bmatrix}
=\Af,
\]
because \(\Shat=\Shat^\top\), \(\What=\What^\top\), and
\(\Bcphat=\Bpchat^\top\). The block locations
define three edge types: person-person, code-code, and person-code. Hence
\(\Af\) is the adjacency matrix of a weighted undirected heterogeneous graph.
\end{proof}

\begin{corollary}[Laplacian positive semidefiniteness]
Let \(D=\diag(d_u)\), where \(d_u=\sum_v A_{uv}\), and define \(L=D-\Af\).
Under the assumptions of the theorem, \(L\) is positive semidefinite.
\end{corollary}

\begin{proof}
For any \(x\in\R^{n+m}\),
\[
x^\top Lx
=\frac12\sum_{u,v}A_{uv}(x_u-x_v)^2\ge 0,
\]
because \(A_{uv}\ge 0\). Thus \(L\succeq 0\).
\end{proof}

\begin{remark}
This result concerns the graph Laplacian, not the adjacency matrix itself. A
nonnegative symmetric adjacency matrix need not be positive semidefinite. For
example, \(\begin{bmatrix}0&1\\1&0\end{bmatrix}\) has eigenvalues \(1\) and
\(-1\).
\end{remark}

\begin{proposition}[Coherent boundary cases]
The SENA formula has the following mathematically coherent limits:
\begin{enumerate}[label=(\roman*)]
\item If \(\gamma=0\), the social and epistemic layers are disconnected.
\item If \(\alpha=\beta=0\), the graph is a pure person-code bipartite graph.
\item If \(B^{PC}=B^{CP}=0\), the formula reduces to a disconnected social block and
epistemic block.
\item If \(\alpha=0\), person-person social edges are suppressed.
\item If \(\beta=0\), code-code epistemic edges are suppressed.
\end{enumerate}
\end{proposition}

\begin{proof}
Each statement follows by substituting the specified zero value into
the fusion matrix and reading the remaining nonzero blocks.
\end{proof}

\begin{theorem}[No direct contradiction with ENA projection]
The claim that SNA and ENA cannot be shown in one model because SNA is not
dimensionally reduced while ENA is dimensionally reduced is false when SENA is
constructed as the SENA fusion matrix and then embedded by a declared graph
embedding operator.
\end{theorem}

\begin{proof}
Let \(\Phi\) be an embedding operator, for example spectral embedding,
Laplacian eigenmaps, multidimensional scaling on graph distances, or a
specified force-directed layout. The displayed coordinates are then
\[
Z=\Phi(\Af)\in\R^{(n+m)\times d}.
\]
The coordinates are derived from the joint object \(\Af\), not by placing raw
SNA nodes into a pre-existing ENA projection. Hence the dimensionality of a
prior ENA display does not prevent a joint SENA embedding. What is invalid is a
different operation: visually overlaying a raw social graph onto an ENA space
and then interpreting all cross-type distances as if they came from one common
model.
\end{proof}

\section{Person-Code-Pair Attribution}

The bridge matrix \(B^{PC}\) answers which person contributed to which code;
\(B^{CP}\) records explicit code-to-target-person direction when observed. ENA
edges, however, are code-code relations. To attribute an ENA edge to persons,
define a person-code-pair tensor or matrix
\[
G_{i,ab}=\text{contribution of }p_i\text{ to the code pair }(c_a,c_b).
\]
A compatible estimator is
\[
G_{i,ab}=\sum_{t=1}^q Y_{it}X_{ta}X_{tb}.
\]
Thus \(B^{PC}\) is appropriate for person-to-code bridge visualization,
\(B^{CP}\) for evidence-backed reverse direction, and \(G\) for claims such as ``person \(i\) supported the
Evidence-Explanation connection.'' The current SENA implementation includes
such a \(G\) block for reporting person-code-pair contribution.

\section{Normalization and Identifiability}

The layer weights \(\alpha,\beta,\gamma\) are interpretable only relative to
the normalization. If \(S,W,B^{PC},B^{CP}\) are rescaled arbitrarily, raw block
magnitudes change their relative influence. The shared \(\gamma\) cannot
compensate arbitrary independent rescaling of the two bridge directions. Thus
the parameters are not intrinsically identifiable without a normalization
convention.

\begin{proposition}[Scale behavior under declared normalization]
For max or Frobenius normalization applied separately to
\(S,W,B^{PC},B^{CP}\), multiplying any raw block by a positive scalar leaves
its normalized block, and therefore \(\Af\), unchanged. Log1p-max normalization
does not have this global scale invariance and must be reported as such.
\end{proposition}

\begin{proof}
For a positively homogeneous norm or maximum, \(N(cM)=cM/(c\,d(M))=N(M)\)
for \(c>0\), where \(d\) is the corresponding divisor. The log1p transform is
not positively homogeneous, so the same cancellation does not hold.
\end{proof}

\begin{corollary}
Any empirical SENA report must disclose the normalization rule before
interpreting \(\alpha,\beta,\gamma\).
\end{corollary}

\section{Directed and Temporal Extensions}

If \(S\ne S^\top\), or if \(B^{CP}\ne(B^{PC})^\top\), then the same block
construction defines a directed heterogeneous graph rather than an undirected
one. The Laplacian result above does not apply without additional directed-graph
choices or declared symmetrization.

The executable directed bridge contract is
\[
\Af=
\begin{bmatrix}
\alpha\Shat & \gamma\Bpchat\\
\gamma\Bcphat & \beta\What
\end{bmatrix},
\]
where \(B^{PC}\) represents person-to-code activation and \(B^{CP}\) represents
code-to-person target, uptake, feedback, or recommendation evidence. Runtime
provenance must identify \texttt{pc-cp-independent} or
\texttt{pc-transpose-fallback}.

For temporal analysis, define
\[
\Af(t)=
\begin{bmatrix}
\alpha\widehat S(t) & \gamma\widehat B^{PC}(t)\\
\gamma\widehat B^{CP}(t) & \beta\widehat W(t)
\end{bmatrix}.
\]
Then temporal change can be studied using
\[
\Delta A(t_1,t_2)=\Af(t_2)-\Af(t_1),
\]
or by distances such as Frobenius distance, spectral distance, graph edit
distance, or permutation-tested group/time contrasts.

\section{Failure Modes and Counterexamples}

\begin{example}[Projection error]
If a researcher places person nodes from an SNA graph onto a previously
computed ENA unit space and interprets person-code distances as distances from
a common model, the interpretation is invalid. The coordinates were not jointly
estimated from \(\Af\).
\end{example}

\begin{example}[Layer dominance]
If \(S\) contains counts in the hundreds while \(W\) and the bridge blocks contain binary
values, then centrality or layout on the unnormalized matrix is dominated by
the social layer. This is a scale artifact, not a substantive social-epistemic
finding.
\end{example}

\begin{example}[Coding unreliability]
If code assignments are unstable, then \(W\), the bridge blocks, and \(G\) inherit that
instability. A formally valid matrix can still be empirically unreliable.
\end{example}

\begin{example}[Bridge provenance error]
If the transpose fallback is reported as independently observed \(B^{CP}\)
evidence, the model overstates directionality. Every report must disclose the
active bridge mode.
\end{example}

\section{Statistical Validation Framework}

A publication-grade SENA study should report:
\begin{enumerate}
\item the data contract: persons, interactions, windows/stanzas, coded
segments, and codebook;
\item exact formulas for \(S,W,B^{PC},B^{CP}\), and, if used, \(G\), including
whether \(B^{CP}\) is independent or the transpose fallback;
\item the normalization rule;
\item sensitivity analyses over \((\alpha,\beta,\gamma)\);
\item the embedding operator \(\Phi\), including random seed if stochastic;
\item coding reliability or human-reviewed AI coding procedures;
\item evidence traceability from each edge to source records;
\item null models or resampling tests for group and temporal claims.
\end{enumerate}

Useful null models include code-label shuffling within learners, social-edge
rewiring with degree preservation, stanza membership shuffling within temporal
bands, group-label permutation, and bootstrap resampling of windows.

\section{Current Tooling Status and Claim Boundary}

As of the late-June 2026 web-tool implementation, SENA has a runnable local
research-pilot workbench with a five-table import contract, local JavaScript
jENA and jSNA runtime handoffs, temporal SENA traces, review-packet exports,
and an enterprise readiness loop for institutional handoff evidence. This
tooling status should not be read as publication-grade empirical validation or
as completed production SaaS cutover. Strong research claims still require real
research datasets, coding reliability evidence, uncertainty or stability
analysis, appropriate null models or resampling tests, and domain expert
review. Until those gates are attached, the correct interpretation of generated
SENA outputs is exploratory and reproducibility-oriented.

\section{Discussion}

The formal analysis supports the mathematical reasonableness of the SENA
formula, but it also narrows the permissible interpretation. The formula does
not magically convert SNA and ENA into one latent space. Rather, it constructs a
joint typed graph. A shared coordinate system becomes meaningful only after a
specified embedding or matrix factorization is applied to that joint graph.

The most defensible theoretical language is therefore:
\begin{quote}
SENA represents collaborative learning as a typed social-epistemic graph whose
supra-adjacency matrix combines person-person social ties, code-code epistemic
co-occurrences, and person-code contribution bridges.
\end{quote}

This wording connects SENA to established mathematics while preserving its
substantive contribution for learning analytics.

\section{Conclusion}

Under explicit dimensional, nonnegativity, symmetry, and normalization
conditions, the SENA fusion matrix
\[
\Af =
\begin{bmatrix}
\alpha\Shat & \gamma\Bpchat\\
\gamma\Bcphat & \beta\What
\end{bmatrix}
\]
is a valid weighted typed adjacency matrix. Its graph Laplacian is positive
semidefinite when \(S\) and \(W\) are symmetric and
\(B^{CP}=(B^{PC})^\top\). Independent target-person evidence activates the
directed bridge extension. Boundary cases reduce to recognizable SNA,
ENA, and bipartite submodels, and joint embeddings are mathematically
permissible when they are computed from \(\Af\) itself. The formula becomes a
rigorous research method only when matrix construction, normalization,
embedding, statistical testing, and evidence traceability are fully specified.

\begin{thebibliography}{99}

\bibitem{bowman2021}
Bowman, D. A., Swiecki, Z., Cai, Z., Eagan, B., Linder, R., Ruis, A. R., ... &
Shaffer, D. W. (2021). The mathematical foundations of epistemic network
analysis. In A. R. Ruis \& S. B. Lee (Eds.), \emph{Advances in Quantitative
Ethnography} (pp. 91--105). Springer.
\url{https://doi.org/10.1007/978-3-030-67788-6_7}.

\bibitem{gasevic2019}
Gasevic, D., Joksimovic, S., Eagan, B. R., \& Shaffer, D. W. (2019). SENS:
Network analytics to combine social and cognitive perspectives of collaborative
learning. \emph{Computers in Human Behavior}, \emph{92}, 562--577.
\url{https://doi.org/10.1016/j.chb.2018.07.003}.

\bibitem{kivela2014}
Kivela, M., Arenas, A., Barthelemy, M., Gleeson, J. P., Moreno, Y., \& Porter,
M. A. (2014). Multilayer networks. \emph{Journal of Complex Networks},
\emph{2}(3), 203--271. \url{https://doi.org/10.1093/comnet/cnu016}.

\bibitem{mucha2010}
Mucha, P. J., Richardson, T., Macon, K., Porter, M. A., \& Onnela, J.-P.
(2010). Community structure in time-dependent, multiscale, and multiplex
networks. \emph{Science}, \emph{328}(5980), 876--878.
\url{https://doi.org/10.1126/science.1184819}.

\bibitem{shaffer2016}
Shaffer, D. W., Collier, W., \& Ruis, A. R. (2016). A tutorial on epistemic
network analysis: Analyzing the structure of connections in cognitive, social,
and interaction data. \emph{Journal of Learning Analytics}, \emph{3}(3),
9--45. \url{https://doi.org/10.18608/jla.2016.33.3}.

\bibitem{siebert2017}
Siebert-Evenstone, A., Arastoopour Irgens, G., Collier, W., Swiecki, Z.,
Ruis, A. R., \& Williamson Shaffer, D. (2017). In search of conversational
grain size: Modelling semantic structure using moving stanza windows.
\emph{Journal of Learning Analytics}, \emph{4}(3), 123--139.
\url{https://doi.org/10.18608/jla.2017.43.7}.

\bibitem{sun2012}
Sun, Y., \& Han, J. (2012). \emph{Mining Heterogeneous Information Networks:
Principles and Methodologies}. Morgan \& Claypool.
\url{https://doi.org/10.2200/S00433ED1V01Y201207DMK005}.

\bibitem{swiecki2020}
Swiecki, Z., \& Shaffer, D. W. (2020). iSENS: An integrated approach to
combining epistemic and social network analyses. In V. Kovanovic, M. Scheffel,
N. Pinkwart, \& K. Verbert (Eds.), \emph{LAK20 Conference Proceedings:
Celebrating 10 Years of LAK: Shaping the Future of the Field} (pp. 305--313).
Association for Computing Machinery.
\url{https://doi.org/10.1145/3375462.3375505}.

\bibitem{wasserman1994}
Wasserman, S., \& Faust, K. (1994). \emph{Social Network Analysis: Methods and
Applications}. Cambridge University Press.

\end{thebibliography}

\end{document}
